\section{Local spectral temperature}
  \label{sec:local-temperature}

  \subsection{Definition of \(\beta(x)\) as a local conjugate field}
    \label{subsec:beta-definition}

    We now introduce a local inverse temperature field \(\beta(x)\) as the conjugate variable to a local energy density
    \(E_\Pi(x)\) associated with the projected sector.
    The operational definition is the local first-law relation
    \begin{equation}
      \delta s_\Pi(x) = \beta(x)^{-1}\,\delta e_\Pi(x),
      \label{eq:local-first-law-density}
    \end{equation}
    where \(s_\Pi(x)\) and \(e_\Pi(x)\) are local densities such that
    \(
    S_\Pi=\int_\Sigma s_\Pi(x)\,d^3x
    \)
    and
    \(
    E_\Pi=\int_\Sigma e_\Pi(x)\,d^3x
    \)
    in a chosen foliation.

    Equation~\eqref{eq:local-first-law-density} is not assumed as microscopic dynamics.
    It is a consistency condition defining \(\beta(x)\) as the local Lagrange multiplier that enforces the
    compatibility between spectral entropy variations and admissible energy variations in the effective description.

  \subsection{Geometric correction to \(\beta(x)^{-1}\)}
    \label{subsec:beta-curvature}

    To extract the curvature correction, we match the local spectral density~\eqref{eq:u-leading} to the local first-law
    structure.
    At fixed spectral scale \(t\), the flat reference defines \(\beta_0^{-1}\) through the leading \(2/t\) term, while the
    first deviation is geometric.
    This yields the decomposition
    \begin{equation}
      \beta(x)^{-1} = \beta_0^{-1} + \frac16 R(x) + O(\nabla^2 R, R^2),
    \end{equation}
    which is~\eqref{eq:beta-curvature}.
    Here the higher terms are controlled by the higher heat-kernel coefficients and by derivative corrections.
    The identification is local and covariant: curvature enters as the first correction to a locally defined spectral
    temperature.
